% Shared preamble for all problems from First_Proof.tex
% Source: https://raw.githubusercontent.com/1stproof/batch-1/refs/heads/main/First_Proof.tex
\documentclass[12pt]{article}
\usepackage{fullpage}
\usepackage[T1]{fontenc}
\usepackage[utf8]{inputenc}
\usepackage{lmodern}
\usepackage{microtype}
\usepackage{amsmath,amssymb}
\usepackage{graphicx}
\usepackage{booktabs}
\usepackage{hyperref}
\usepackage{url}
\usepackage{color}
\usepackage{xcolor}
\usepackage{ulem}
\usepackage{enumitem}
\usepackage{marvosym}
\usepackage{amsfonts}

\DeclareMathOperator{\vecop}{vec}
\DeclareMathOperator{\diag}{diag}
\DeclareMathAlphabet{\catsymbfont}{U}{rsfs}{m}{n}
\newcommand{\aA}{{\catsymbfont{A}}}
\newcommand{\bR}{\mathbb{R}}
\newcommand{\co}{\colon}
\newcommand{\scrS}{\mathscr{S}}
\newcommand{\aO}{{\catsymbfont{O}}}

\usepackage{amsthm}
\newtheorem{lemma}{Lemma}
\newtheorem{theorem}{Theorem}
\newtheorem{remark}{Remark}

\title{Problem 10 (Lemma L1): The normal-equation operator is symmetric positive (semi)definite}
\date{}

\begin{document}
\maketitle

\noindent
This note proves the spectral properties of the mode-$k$ normal-equation
operator that are the prerequisite for using the conjugate gradient (CG)
method. Throughout, with the notation of the problem statement, set
\begin{equation}\label{eq:Adef}
  \aA \;=\; (Z\otimes K)^{T} S\, S^{T} (Z\otimes K)\;+\;\lambda\,(I_r\otimes K)
  \;\in\;\bR^{nr\times nr},
\end{equation}
and let
\begin{equation}\label{eq:rhs}
  b \;=\; (I_r\otimes K)\,\vecop(B),\qquad B = TZ,
\end{equation}
so that the subproblem is $\aA\,\vecop(W)=b$. We always assume, as in the
problem statement, that $K=K^{T}\in\bR^{n\times n}$ is positive semidefinite
($K\succeq0$, the RKHS kernel matrix), that $S\in\bR^{N\times q}$ is a column
subset of $I_N$ (hence $S^{T}S=I_q$ and $SS^{T}$ is an orthogonal projector),
and that $\lambda\ge0$.

\bigskip

\begin{lemma}[Symmetry and positive semidefiniteness]\label{lem:psd}
For every $K\succeq0$ and every $\lambda\ge0$ the operator $\aA$ in
\eqref{eq:Adef} is symmetric and positive semidefinite, i.e.
$\aA=\aA^{T}$ and $x^{T}\aA\,x\ge0$ for all $x\in\bR^{nr}$.
\end{lemma}

\begin{proof}
\emph{Symmetry.} Write $G:=(Z\otimes K)^{T} S\,S^{T} (Z\otimes K)$ and
$R:=\lambda(I_r\otimes K)$, so $\aA=G+R$.
For any matrix $M:=S^{T}(Z\otimes K)\in\bR^{q\times nr}$ we have
$G=M^{T}M$, which is symmetric because $(M^{T}M)^{T}=M^{T}M$.
For $R$, the transpose of a Kronecker product satisfies
$(I_r\otimes K)^{T}=I_r^{T}\otimes K^{T}=I_r\otimes K$ because $I_r$ and $K$
are symmetric; multiplying by the scalar $\lambda$ preserves symmetry, so
$R^{T}=R$. Hence $\aA^{T}=G^{T}+R^{T}=G+R=\aA$.

\smallskip
\emph{Positive semidefiniteness.} Fix $x\in\bR^{nr}$. Then
\begin{equation}\label{eq:quadform}
  x^{T}\aA\,x
  = x^{T}M^{T}M x + \lambda\, x^{T}(I_r\otimes K)x
  = \underbrace{\|Mx\|_2^{2}}_{\ge0}
   \;+\; \lambda\, x^{T}(I_r\otimes K)x .
\end{equation}
The first term is a squared Euclidean norm, hence $\ge0$. For the second term,
since $K\succeq0$ there is a factorization $K=L L^{T}$ (e.g.\ a Cholesky or
symmetric square-root factor $L\in\bR^{n\times n}$). Using the Kronecker
identity $I_r\otimes K=I_r\otimes(LL^{T})=(I_r\otimes L)(I_r\otimes L)^{T}$,
\[
  x^{T}(I_r\otimes K)x
  = x^{T}(I_r\otimes L)(I_r\otimes L)^{T}x
  = \bigl\|(I_r\otimes L)^{T}x\bigr\|_2^{2}\;\ge\;0 .
\]
Because $\lambda\ge0$, the second term in \eqref{eq:quadform} is $\ge0$, and
therefore $x^{T}\aA x\ge0$. Thus $\aA\succeq0$.
\end{proof}

\begin{lemma}[Strict positive definiteness]\label{lem:spd}
Assume in addition that $K\succ0$ (i.e.\ $K$ is nonsingular) and $\lambda>0$.
Then $\aA$ is symmetric \emph{positive definite}: $x^{T}\aA x>0$ for all
$x\neq0$, equivalently $\aA$ is nonsingular. These two conditions are also
\emph{necessary} in general: if $\lambda=0$ then $\aA$ is singular, and if
$K$ is singular then $\aA$ is singular.
\end{lemma}

\begin{proof}
\emph{Sufficiency.} Symmetry holds by Lemma~\ref{lem:psd}. Let $x\neq0$.
By \eqref{eq:quadform},
$x^{T}\aA x \ge \lambda\,x^{T}(I_r\otimes K)x$. Since $K\succ0$, its
eigenvalues are bounded below by $\mu_{\min}(K)>0$, and the eigenvalues of the
Kronecker product $I_r\otimes K$ are exactly the eigenvalues of $K$, each with
multiplicity $r$ (a standard fact: if $K v=\mu v$ then
$(I_r\otimes K)(e_j\otimes v)=\mu\,(e_j\otimes v)$ for each standard basis
vector $e_j$, giving $nr$ orthonormal eigenpairs in total). Hence
$I_r\otimes K\succeq \mu_{\min}(K)\,I_{nr}$ and
\[
  x^{T}\aA x \;\ge\; \lambda\,\mu_{\min}(K)\,\|x\|_2^{2} \;>\;0 ,
\]
because $\lambda>0$, $\mu_{\min}(K)>0$, and $x\neq0$. A symmetric operator
with strictly positive quadratic form on all nonzero vectors is positive
definite and in particular nonsingular.

\smallskip
\emph{Necessity of $\lambda>0$.} If $\lambda=0$, then
$\aA = (Z\otimes K)^{T}SS^{T}(Z\otimes K) = M^{T}M$ with
$M=S^{T}(Z\otimes K)\in\bR^{q\times nr}$. Since $M$ has only $q$ rows and
$\aA$ has order $nr$, whenever $q<nr$ the rank of $\aA$ is at most $q<nr$, so
$\aA$ is singular. (Under the standing assumption $n,r<q\ll N$ one typically
has $q\ge n,r$ individually, but not $q\ge nr$; in the regime of interest
$q\ll N=nM$ and $nr$ can exceed $q$, so the $\lambda=0$ operator is singular.)
Thus $\lambda>0$ is required for guaranteed nonsingularity.

\smallskip
\emph{Necessity of $K\succ0$.} Suppose $K$ is singular, so there is
$v\neq0$ with $Kv=0$. Pick any $e_j$ ($j\in\{1,\dots,r\}$) and set
$x:=e_j\otimes v\neq0$. Using the mixed-product rule
$(A\otimes B)(C\otimes D)=(AC)\otimes(BD)$,
\[
  (Z\otimes K)x=(Z\otimes K)(e_j\otimes v)=(Ze_j)\otimes(Kv)=(Ze_j)\otimes 0=0,
  \qquad
  (I_r\otimes K)x=(I_r e_j)\otimes(Kv)=e_j\otimes 0=0 .
\]
Hence $\aA x=(Z\otimes K)^{T}SS^{T}\cdot 0+\lambda(I_r\otimes K)x=0$ for
\emph{any} $\lambda$, so $x\in\ker\aA$ and $\aA$ is singular.
\end{proof}

\begin{remark}[Why this matters for CG]\label{rem:cg}
The (unpreconditioned) CG method is guaranteed to converge to the unique
solution when applied to a symmetric positive definite system. By
Lemma~\ref{lem:spd}, the hypotheses
\[
  \boxed{\;K\succ0\quad\text{and}\quad\lambda>0\;}
\]
are exactly what guarantees $\aA\succ0$, hence the well-posedness of CG.
Many RKHS kernel matrices are strictly positive definite (e.g.\ the Gaussian
kernel on distinct sample points), in which case this hypothesis holds
verbatim. The remaining (boundary) case $K\succeq0$ singular is treated next.
\end{remark}

\begin{lemma}[Consistency and CG in the singular--$K$ boundary case]
\label{lem:singular}
Let $K\succeq0$ be (possibly) singular and $\lambda>0$. Then:
\begin{enumerate}[label=\textup{(\roman*)}]
\item $\ker\aA=\ker(I_r\otimes K)$ and hence
      $\operatorname{range}(\aA)=\operatorname{range}(I_r\otimes K)$;
\item the right-hand side satisfies
      $b=(I_r\otimes K)\vecop(B)\in\operatorname{range}(I_r\otimes K)
      =\operatorname{range}(\aA)$, so the system $\aA\,w=b$ is
      \emph{consistent} (solvable);
\item CG started from $w_0=0$ produces iterates $w_m\in\operatorname{range}(\aA)$
      and converges, in exact arithmetic, to the unique minimum-norm solution
      $w_\star\in\operatorname{range}(\aA)$ in at most $\operatorname{rank}(\aA)$
      iterations.
\end{enumerate}
In particular CG remains applicable even when $K$ is only semidefinite, with
the convention that it solves the system on $\operatorname{range}(\aA)$.
\end{lemma}

\begin{proof}
(i) Write $\aA=M^{T}M+\lambda(I_r\otimes K)$ as in
Lemma~\ref{lem:psd} with $M=S^{T}(Z\otimes K)$, and note both summands are
PSD. For any PSD matrices $P,Q$ one has
$\ker(P+Q)=\ker P\cap\ker Q$: indeed if $(P+Q)x=0$ then
$0=x^{T}(P+Q)x=x^{T}Px+x^{T}Qx$ with both terms $\ge0$, forcing
$x^{T}Px=x^{T}Qx=0$, and PSD-ness gives $Px=Qx=0$; the reverse inclusion is
immediate. Applying this with $P=M^{T}M$, $Q=\lambda(I_r\otimes K)$,
\[
  \ker\aA=\ker(M^{T}M)\cap\ker\!\bigl(\lambda(I_r\otimes K)\bigr)
        =\ker M\cap\ker(I_r\otimes K).
\]
But $\ker(I_r\otimes K)\subseteq\ker M$: using
$(Z\otimes K)=(Z\otimes I_n)(I_r\otimes K)$ via the mixed-product rule, any
$x$ with $(I_r\otimes K)x=0$ satisfies $(Z\otimes K)x=0$ and hence
$Mx=S^{T}(Z\otimes K)x=0$. Therefore
$\ker\aA=\ker(I_r\otimes K)$. Since $\aA$ and $I_r\otimes K$ are symmetric,
their ranges are the orthogonal complements of their kernels, so equal kernels
give equal ranges:
$\operatorname{range}(\aA)=\ker(\aA)^{\perp}=\ker(I_r\otimes K)^{\perp}
=\operatorname{range}(I_r\otimes K)$.

\smallskip
(ii) By definition $b=(I_r\otimes K)\vecop(B)$ is in the column space of
$I_r\otimes K$, i.e.\ $b\in\operatorname{range}(I_r\otimes K)$. By (i) this
equals $\operatorname{range}(\aA)$, so there exists $w$ with $\aA w=b$; the
system is consistent. (This consistency is structural: the right-hand side is
itself produced by the operator $I_r\otimes K$ that defines the kernel of
$\aA$, so no component of $b$ ever falls in $\ker\aA$.)

\smallskip
(iii) Since $\aA$ is symmetric PSD and the system is consistent, CG is the
minimization of the convex quadratic
$\phi(w)=\tfrac12 w^{T}\aA w-b^{T}w$ over Krylov subspaces. With $w_0=0$ the
initial residual is $r_0=b\in\operatorname{range}(\aA)$, and the Krylov
subspace $\operatorname{span}\{r_0,\aA r_0,\dots,\aA^{m-1}r_0\}\subseteq
\operatorname{range}(\aA)$ because $\operatorname{range}(\aA)$ is
$\aA$-invariant. Hence every iterate $w_m$ lies in
$\operatorname{range}(\aA)$, the iteration never ``sees'' $\ker\aA$, and CG
behaves exactly as it would on the invertible restriction
$\aA|_{\operatorname{range}(\aA)}$, which is symmetric positive definite of
order $\operatorname{rank}(\aA)$. Therefore the residuals decrease and, in
exact arithmetic, $r_m=0$ after at most $\operatorname{rank}(\aA)$ distinct
eigenvalues' worth of iterations, with limit the unique solution in
$\operatorname{range}(\aA)$, namely the minimum-norm solution $w_\star$.
\end{proof}

\begin{remark}[Practical safeguard]\label{rem:reg}
When $K\succeq0$ is (numerically) singular and one prefers a strictly SPD
system, replace $K$ by $K_\varepsilon:=K+\varepsilon I_n$ with a small
$\varepsilon>0$. Then $K_\varepsilon\succ0$, so Lemma~\ref{lem:spd} applies and
$\aA_\varepsilon=(Z\otimes K_\varepsilon)^{T}SS^{T}(Z\otimes K_\varepsilon)
+\lambda(I_r\otimes K_\varepsilon)\succ0$ for every $\lambda>0$; CG then
converges to the unique solution of $\aA_\varepsilon w=b_\varepsilon$, where
$b_\varepsilon=(I_r\otimes K_\varepsilon)\vecop(B)$. As $\varepsilon\to0^{+}$
this regularized solution converges to the minimum-norm solution
$w_\star$ of the original consistent system of Lemma~\ref{lem:singular},
recovering the boundary case continuously. Crucially, both the matrix--vector
products and the preconditioner used elsewhere in the algorithm depend only on
the factor/kernel data, so this shift adds no asymptotic cost.
\end{remark}

\begin{remark}[The limit $\lambda\to0$]\label{rem:lam0}
For $\lambda=0$ the operator reduces to $\aA=M^{T}M$ with $M\in\bR^{q\times nr}$,
which is singular whenever $q<nr$ (as in the regime $q\ll N$ of interest), and
the consistency argument of Lemma~\ref{lem:singular}(ii) no longer applies
because the right-hand side $b=(I_r\otimes K)\vecop(B)$ need not lie in
$\operatorname{range}(M^{T}M)=\operatorname{range}(M^{T})$. Hence $\lambda>0$ is
an essential well-posedness hypothesis, not merely a convenience: it both
provides strict positive definiteness (with $K\succ0$) and, in the singular-$K$
case, aligns $\operatorname{range}(\aA)$ with $\operatorname{range}(I_r\otimes K)$
so that the system stays consistent.
\end{remark}

\end{document}
